%% kapa: exam=1 round=34 year=2023 period=1 form=A subject=civil_law subject_ko=민법 seq=35 number=35
%% source: Q-Net 감정평가사 기출문제 (cst00309) artlSeq=5211789 file=raw/1008/5211789/2023년도 제34회 감정평가사 1차시험 1교시.pdf
%% conversion: kordoc 4.24.51 -> tools/gen1.py (auto); review=proofread-sonnet; flags=none
\begin{question}{35}
전세권에 관한 설명으로 옳은 것은? (다툼이 있으면 판례에 따름)
\begin{choices}
  \item 건물 일부의 전세권자는 나머지 건물 부분에 대해서도 경매신청권이 있다.
  \item 전세권 설정계약의 당사자는 전세권의 사용·수익권능을 배제하고 채권담보만을 위한 전세권을 설정할 수 있다.
  \item 전세권설정시 전세금 지급은 전세권 성립의 요소이다.
  \item 전세권자는 특별한 사정이 없는 한 전세권의 존속기간 내에서 전세목적물을 타인에게 전전세 할 수 없다.
  \item 전세권이 소멸된 경우, 전세권자의 전세목적물의 인도는 전세금의 반환보다 선이행되어야 한다.
\end{choices}
\end{question}
\answer{3}
